%==============================================================================
\section{Discussão}
\label{sec:discussion}

% --- corresponde a 01_Project_Overview.md §6 (contribuição) e §7 (ressalvas) ---

\subsection{O que a teoria da informação acrescenta sobre $R^2$}
Uma métrica de ruído correta e transferível: a informação mútua não assume
uma forma funcional, é válida onde a física é não-linear e é expressa em
bits, ao contrário da fração de variância. Sobre a série bruta --- o
pareamento correto com o $R^2$ publicado do trabalho anterior, também
medido sobre série bruta --- a dependência medida excede o equivalente
linear-$R^2$ em $\valMIexcessPctBinned$--$\valMIexcessPctKSG\,\%$. Esse
excesso é real (sobrevive ao piso de viés do estimador,
\cref{sec:res-mi}), mas \textbf{não é, por si só, evidência de
não-linearidade instantânea na relação reversível $\Tdie,\Vcore\to\slk$}: a
mesma comparação sobre as séries detrendizadas mostra um \emph{déficit}, não
um excesso ($\valMIexcessPctDetrended\,\%$, \cref{sec:res-mi}), de modo que
o excesso bruto está concentrado na tendência lenta que envelhecimento e
$\Tdie,\Vcore$ compartilham ao longo da campanha, não na física
Arrhenius/BTI instantânea que motivou este projeto. A contribuição
metodológica --- MI em vez de $R^2$ como a métrica correta e transferível,
válida mesmo quando a forma da dependência é desconhecida --- permanece de
pé; a interpretação causal específica (``a MI excede o $R^2$ porque a
física subjacente é não-linear'') não é sustentada por este resultado e é
retirada. Isso também enfraquece a extrapolação entre regimes térmicos: mesmo a
leitura mais cautelosa --- que a comparação anterior, baseada em $R^2$,
entre bang-bang e PID poderia subestimar a diferença real entre eles ---
perde parte de sua base, já que o excesso não-linear que a sustentaria não
se mostrou robusto ao \emph{detrend}; recalcular ambas as bancadas lado a
lado em bits exigiria, de qualquer forma, uma execução pareada no mesmo
dispositivo em ambos os regimes, fora do escopo deste estudo de regime
único (\cref{sec:discussion}, ameaças à validade).

\subsection{Uma definição principiada de qualidade de \emph{burn-in}}
Reenquadrar a qualidade de \emph{burn-in} como a minimização de
$\MI{\slk}{\Tdie,\Vcore}$ dá à afirmação qualitativa ``PID é um ambiente de
validação mais limpo'' uma unidade rigorosa e um alvo de otimização: cada
bit que o ambiente escreve na leitura de \emph{slack} é um bit de
informação de envelhecimento perdido. Enunciamos isso como um princípio, em
vez de demonstrá-lo dentro de um único dispositivo; o suporte empírico mais
próximo aqui é a lacuna distribucional entre as duas bancadas
(\cref{sec:res-kl}), medida em um dispositivo de teste separado.

\subsection{Uma visão por comprimento de descrição da escolha linear-vs-não-linear}
O critério MDL da \cref{sec:th-mdl}, que seleciona a correção PVT linear
(\cref{sec:res-correction}), não é uma penalidade \emph{ad hoc} sobre
parâmetros extras: é um substituto computável para a intuição de
complexidade de Kolmogorov de que o melhor modelo é a \emph{descrição mais
curta} que ainda explica os dados (\cref{eq:kolmogorov}). Enquadrado dessa
forma, o resultado não é meramente ``a correção não-linear não vale a pena
aqui'' --- é que, neste único e estreito ponto de operação
(\valDutTempMean~$^\circ$C, \valDutVoltMean~V), os dados ainda não contêm
informação de curvatura suficiente para justificar pagar pelos parâmetros
extras do modelo Arrhenius/super-linear, mesmo que a física subjacente seja
genuinamente não-linear (\cref{sec:background-r2}). Isso é, na verdade,
\textbf{a mesma conclusão} alcançada por uma rota independente: o excesso
de MI sobre o equivalente Gaussiano linear, que motivou este projeto, não
sobrevive quando medido sobre o acoplamento reversível detrendizado
(\cref{sec:res-mi}). MDL (nesta seção) e a checagem de MI detrendizada
(\cref{sec:res-mi}) perguntam coisas diferentes --- se uma forma paramétrica
não-linear específica se paga na correção, e se o acoplamento reversível
\emph{instantâneo} contém excesso de dependência não-linear, respectivamente
--- mas os dois, de forma independente, dizem que este conjunto de dados de
único ponto de operação não contém evidência de não-linearidade
\emph{instantânea} entre $\Tdie,\Vcore$ e $\slk$, mesmo que a física
subjacente seja genuinamente não-linear em uma faixa mais ampla de operação.
Uma campanha de faixa mais larga ou com múltiplos pontos de ajuste, cobrindo
curvatura Arrhenius/BTI suficiente para ser compressível pelo modelo
não-linear, é a condição natural sob a qual ambos os vereditos se
inverteriam.

\subsection{Um limite de resolução para sensores de SLM}
Os resultados de capacidade e BSC colocam números concretos sobre o que o
sensor consegue fazer: \valCapacityStates{} estados de envelhecimento
distinguíveis ao longo da campanha, um alarme diário
($\valBSCwindow$~h) de envelhecimento de $\Cbsc=\valBSCcap$~bits, e um piso
de detecção a $3\sigma$ de \valMinDetectRate~m-cnt/h atingido após
$\sim$\valMinObsTime~h. Isso informa a um projetista de SLM quão finamente o
sensor resolve o envelhecimento e quão confiavelmente ele pode disparar um
alarme sob ruído realista --- grandezas diretamente consumidas por
escalonamento adaptativo de tensão e RUL, e que a métrica residual-$\sigma$
anterior deixava implícitas. Sua consistência mútua (a confiabilidade do
BSC por janela cruza próximo ao tempo mínimo de observação do CRLB) é
evidência de que o modelo de canal é coerente, e não três heurísticas
desconectadas.

\subsection{Estimação sequencial e RUL como distribuição}
O tratamento Wiener/Kalman da \cref{sec:res-wiener} deliberadamente não é um
substituto do limite de Cram\'er--Rao/Bayesiano da \cref{sec:res-rul}, mas
uma verificação cruzada a partir de uma rota de estimação independente (um
modelo dinâmico e filtrado, em vez de uma única regressão sobre toda a
série) que, por sinal, concorda no sinal e na ordem de grandeza da taxa de
envelhecimento. Seu valor agregado é estrutural, não numérico: transforma
uma única taxa pontual com barra de erro em (i) uma estimativa de
trajetória denoised, continuamente atualizada, que o operador poderia, em
princípio, consumir on-line, e (ii) uma \emph{distribuição} completa de RUL,
em vez de um limite, expondo a assimetria à direita que uma barra de erro
Gaussiana esconde. O fato de que o MLE do ruído de taxa e o critério MDL da
\cref{sec:res-correction} \emph{independentemente} preferem, ambos, o modelo
mais simples disponível (deriva constante; correção linear) neste único
ponto de operação é, por si só, um achado: dois princípios de estimação
diferentes, aplicados a duas perguntas diferentes, concordam que esta
campanha ainda não contém a curvatura necessária para justificar um modelo
mais flexível --- exatamente o tipo de validação cruzada que este projeto
usou ao longo de todo o trabalho (KSG vs.\ MI com \emph{binning},
frequentista vs.\ RUL Bayesiano, ICA vs.\ correção por regressão).

\subsection{Ameaças à validade}
O excesso de informação mútua sobre o equivalente Gaussiano-linear, medido
na série bruta (\cref{sec:res-mi}), não sobrevive à remoção da tendência de
envelhecimento compartilhada e não deve ser lido como demonstração de
não-linearidade instantânea na física PVT reversível --- esse é o achado
mais importante desta seção de ameaças à validade, não uma nota de rodapé,
e é discutido em detalhe nas \cref{sec:res-mi,sec:discussion}. A análise do
dispositivo de envelhecimento é de um único dispositivo e um
único regime térmico, o que limita a separabilidade do envelhecimento
permanente. A comparação de bancadas da \cref{sec:res-kl} é em um
dispositivo de teste \emph{diferente}, de modo que sua divergência KL/JS
fala sobre a fidelidade de bancada, e não sobre o comportamento do
dispositivo de envelhecimento, e os dois nunca são misturados. Estimativas
de MI em dados autocorrelacionados são propensas a viés, apesar da grande
contagem bruta de amostras, mitigado (não eliminado) por $\neff$, pelo piso
nulo \emph{surrogate} da \cref{sec:res-mi} e pela comparação entre
estimadores independentes (\emph{binning} vs.\ KSG), cuja diferença
sistemática é reportada como faixa de incerteza, não escondida; a
expressão de capacidade é uma heurística de resolução
efetiva; a separação por ICA é corroborativa, não uma desmistura limpa. O
modelo Wiener/Kalman da \cref{sec:res-wiener} assume ruído de observação por
média de bloco constante e um processo latente genuinamente difusivo
(Wiener); a distribuição de RUL que ele produz é reportada em limiares
auto-referenciados, já calculados, precisamente porque não existe um limiar
de falha calibrado de forma independente para este sensor, de modo que é
uma demonstração de \emph{metodologia}, não uma predição operacional de
vida útil, no mesmo espírito do limite de Cram\'er--Rao da
\cref{sec:res-rul}.
